\documentclass[11pt,UTF8,fontset=fandol]{ctexart}
\usepackage{amsmath,amssymb,amsfonts}
\usepackage{graphicx,float,caption,subcaption}
\usepackage{booktabs,array,multirow,makecell}
\usepackage{geometry,enumitem,xcolor}
\usepackage[hidelinks]{hyperref}
\usepackage{url}
\geometry{margin=2.3cm}
\graphicspath{{/Users/yye/project/论文/assets/llaga/}}
\captionsetup{font=small,labelfont=bf}
\setlength{\parskip}{0.4em}
\hypersetup{pdftitle={LLaGA 中文精读导读},pdfauthor={Chen et al.}}

% 小白知识卡片：上下加线 + 粗体标签，稳健不挑内容
\newenvironment{tip}[1]{%
  \par\smallskip\noindent\rule{\linewidth}{0.4pt}\par\nobreak
  \noindent\textbf{【小白补丁】#1}\par\nobreak\noindent\ignorespaces
}{%
  \par\noindent\rule{\linewidth}{0.4pt}\par\smallskip
}

\title{LLaGA 中文精读导读：把一个节点“翻译”成大模型听得懂的“外语词”}
\author{原文：Runjin Chen, Tong Zhao, Ajay Jaiswal, Neil Shah, Zhangyang Wang（ICML 2024）}
\date{}

\begin{document}
\maketitle
\begin{center}\small\textit{本导读为面向初学者的中文学习笔记：含要点翻译、公式精讲与图表解读。}\end{center}

\section*{论文卡片}
\begin{center}
\begin{tabular}{@{}ll@{}}
\toprule
\textbf{标题} & LLaGA: Large Language and Graph Assistant \\
            & （大语言与图助手） \\
\textbf{作者} & Runjin Chen, Tong Zhao, Ajay Jaiswal, Neil Shah, Zhangyang Wang \\
\textbf{机构} & 德州大学奥斯汀分校（UT Austin）、Snap Inc. \\
\textbf{会议} & ICML 2024（国际机器学习大会） \\
\textbf{arXiv} & 2402.08170 \\
\textbf{代码} & \url{https://github.com/VITA-Group/LLaGA} \\
\textbf{方向} & 图（Graph）$\times$ 大模型 $\times$ 后训练（投影器对齐 / 表示对齐） \\
\textbf{一句话} & 把一个节点和它的邻居重排成“结构感知 token 序列”，用一个轻量\\
              & 投影器映射进 LLM 词向量空间当“外语词”，冻结 LLM、只训投影器，\\
              & 一个统一模型同时做节点分类 / 链接预测 / 节点描述，还能零样本迁移。 \\
\bottomrule
\end{tabular}
\end{center}

\section{一、这篇论文在解决什么问题}

现实世界里大量数据天生是“图”：社交网络、生物网络、推荐系统、引文网络……
它们由\textbf{节点}（实体）和\textbf{边}（关系）组成。专门处理图的工具叫
\textbf{图神经网络（Graph Neural Network, GNN）}，它靠“消息传递 + 聚合”把每个节点
变成一个向量，在图任务上很强。

另一边，\textbf{大语言模型（Large Language Model, LLM，如 GPT-4、LLaMA）}最近大放异彩——
\textbf{一个模型就能干很多种任务}，还会用自然语言解释自己的答案。于是一个自然的念头冒出来：
能不能让 LLM 也来处理图？但作者指出，已有做法各有硬伤：

\begin{enumerate}[leftmargin=2em]
  \item \textbf{GNN 不擅长多任务}：GNN 通常只为“某一个任务”训练，换个任务往往就要换“任务专用的脑袋
        （分类头）”重新调，做不到一个模型通吃。
  \item \textbf{把图“说”成文字效果差}：有人尝试把图结构直接用自然语言描述（比如“节点 A 连着 B 和 C……”）
        喂给 LLM。但这种描述\textbf{又啰嗦又不直观}，LLM 在基础图任务上反而表现很差。
  \item \textbf{指令微调会“偏科”}：像 InstructGLM 那样为某个图任务专门微调 LLM，虽然该任务变好了，
        却\textbf{牺牲了 LLM 的通用性}，在别的任务上变弱。
  \item \textbf{预训练图编码器未必合身}：像 GraphGPT 那样先用一个预训练好的图 Transformer 抽结构信息，
        但这个“现成编码器”未必能为每个下游任务蒸馏出恰好需要的结构特征。
\end{enumerate}

\textbf{本文要回答的核心问题}：怎样设计一个框架，能\textbf{跨多种任务、多种领域}有效地把图的结构信息
编码出来，让 LLM 看得懂，\textbf{同时又不破坏 LLM 原本的通用能力}？这就是 LLaGA。

\section{二、摘要中译}

\begin{quote}
图神经网络（GNN）推动了图结构数据分析的发展。近来，以 GPT-4 为代表的大语言模型（LLM）
开启了深度学习的新纪元；然而，由于“把图结构翻译成语言”本身就很困难，将 LLM 应用到图数据上
仍面临独特挑战。为此，我们提出\textbf{大语言与图助手（Large Language and Graph Assistant，LLaGA）}，
一个能有效整合 LLM 能力来处理复杂图结构数据的创新模型。LLaGA 在\textbf{保留 LLM 通用性}的同时，
把图数据改造成与 LLM 输入兼容的格式。具体做法是：把图中的节点\textbf{重新组织成结构感知的序列}，
再通过一个\textbf{通用投影器（versatile projector）}把它们映射进 token 词向量空间。LLaGA 在\textbf{通用性、
泛化性与可解释性}三方面都很出色：它能在不同数据集和任务上始终表现良好，能把能力扩展到\textbf{未见过的
数据集或任务}，还能为图提供解释。我们在主流图基准上的大量实验表明，LLaGA 用\textbf{一个模型}就在
四个数据集、三种任务上取得出色性能，在监督与零样本两种场景下都超过了当前最强的图模型。
代码见 \url{https://github.com/VITA-Group/LLaGA}。
\end{quote}

\section{三、三大核心贡献（论文称为三个特点）}
\begin{enumerate}[leftmargin=2em]
  \item \textbf{通用性（Versatility）}：用一个简单却普适的方法编码图的结构细节，再用\textbf{单个通用投影器}
        把图空间和 token 空间对齐。这个投影器能跨多个数据集处理多种图任务，\textbf{不需要任务专用的改动}——
        而且效果甚至能超过那些为单一任务量身定做的专用图模型。
  \item \textbf{泛化性（Generalizability）}：因为图与 token 空间被全面对齐了，LLaGA 不仅在训练见过的数据/任务上强，
        还能\textbf{不做任何额外微调}就泛化到\emph{从没见过}的数据集和任务（零样本）。
  \item \textbf{可解释性（Interpretability）}：LLaGA 能对节点嵌入给出\textbf{详细的文字解释}（节点描述任务），
        让人能看懂它“为什么这么判断”。
\end{enumerate}

\section{四、预备知识（小白先看这里）}

\begin{tip}{什么是“文本属性图”（text-attributed graph, TAG）}
就是\textbf{节点上带文字}的图。比如引文网络里，每个节点是一篇论文，它的“文字属性”就是论文的标题+摘要；
电商网络里每个节点是一件商品，文字属性是商品描述。本文主要针对这种图。有了文字，就能用语言模型
把它编码成向量。
\end{tip}

\begin{tip}{什么是“token 嵌入空间”和“投影器（projector）”}
LLM 内部不直接读文字，而是先把每个词（token）查成一个向量，这些向量所在的空间就叫\textbf{token 嵌入空间
（词向量空间）}。本文的“投影器”是一个很小的神经网络（一个 MLP），作用是把“图节点的向量”\textbf{搬到}
这个词向量空间里，让它伪装成 LLM 眼中的一个“词”。这正是 LLaVA 等“多模态 LLM”把图像塞进 LLM 的同款套路，
只不过这里塞的是“图”。
\end{tip}

\begin{tip}{什么是“冻结 LLM、只训投影器”（这就是“后训练”的体现）}
传统微调要改 LLM 的全部参数，又慢又贵，还容易让它“忘掉”原有的语言本领。本文反其道而行：
\textbf{把 LLM 整个冻住、一个参数都不改}，只训练那个小投影器。代价极小，却能把“图的结构信息”对齐进 LLM。
\end{tip}

形式化记号（看不懂可跳过，用到再回看）：一张图记作 $G=(V,E,X)$，其中 $V$ 是节点集、$E$ 是边集、
$X$ 是节点的属性信息。每个节点 $v_i\in V$ 带一个属性特征 $x_i\in X$（本文里是\textbf{文字}）。
另外用 $N_v^k$ 表示节点 $v$ 的\textbf{第 $k$ 跳邻居集合}（即“走 $k$ 步能到的邻居”）。
用一个现成的文本编码模型 $\phi$（如 SBERT、RoBERTa、SimTeG）把节点文字编码成向量 $\phi(x_v)$——
这一步把“文字”变成“能算相似度的数字向量”，是后面一切的基础。

\section{五、数据集与任务：一个模型，三种任务，四个数据集}

\textbf{四个数据集}（详细统计见原文 Table 7）。引文网络里节点是论文、边是共同引用；电商网络里节点是商品、
边是共同购买：

\begin{table}[H]\centering
\caption{LLaGA 用到的四个数据集（数据源自原文 Table 7）。}
\begin{tabular}{lcrr}
\toprule
\textbf{数据集} & \textbf{领域} & \textbf{节点数} & \textbf{边数} \\
\midrule
Cora     & 引文网络（citation）     & 2{,}708     & 5{,}429 \\
Pubmed   & 引文网络（citation）     & 19{,}717    & 44{,}338 \\
ogbn-Arxiv     & 引文网络（citation）     & 169{,}343   & 1{,}166{,}243 \\
ogbn-Products  & 电商网络（e-commerce）   & 2{,}449{,}029 & 61{,}859{,}140 \\
\bottomrule
\end{tabular}
\end{table}

可见数据集规模从几千节点（Cora）一路跨到两百多万节点（Products），稀疏程度也各不相同，能很好地考验方法。

\textbf{三种任务}（一个模型全包）：
\begin{itemize}[leftmargin=2em]
  \item \textbf{节点分类（node classification）}：判断一个节点属于哪一类（论文属于哪个研究领域、商品属于哪个品类）。
        评测指标是\textbf{准确率（Accuracy）}。
  \item \textbf{链接预测（link prediction）}：判断两个节点之间是否应该有一条边。指标也是准确率。
  \item \textbf{节点描述（node description）}：让模型\textbf{用一段话描述某个节点}（“这个节点代表一篇关于……的论文”）。
        这是本文比较新颖的任务，用来把“节点嵌入”对齐到“具体的描述文字”，从而带来可解释性。
        指标是 \textbf{Sbert Score}（生成描述与真值描述的语义相似度）和 \textbf{Description Label Accuracy}
        （从描述里反推出的类别是否正确）。
\end{itemize}

\begin{tip}{统一成“问答（QA）格式”是关键巧思}
传统 GNN 做多任务，要给每个任务配一个“专用脑袋”和“专用损失函数”。LLaGA 不这样——它把\textbf{所有任务都写成
同一种“问答对话”格式}：给一段提示词，里面塞入“节点序列”，让 LLM 像聊天一样把答案生成出来。
于是不再需要任务专用的损失或分类头，一个统一框架就能通吃三种任务。
\end{tip}

\section{六、方法逐节精讲}

LLaGA 的核心只有两步（见图~\ref{fig:arch}）：\textbf{第一步}把“图结构”变成“节点序列”（这一步完全
\emph{不含参数}，纯靠规则重排）；\textbf{第二步}用投影器把“节点序列”变成“token 嵌入序列”塞进\emph{冻结的}
LLM。整个训练里，\textbf{唯一要学的参数就是那个投影器}。

\begin{figure}[H]\centering
  \includegraphics[width=0.98\linewidth]{cand_p03_k00.pdf}
  \caption{LLaGA 总体框架与提示词设计范式（原文 Figure 1，图内英文为原文保留）。
  \textbf{Step 1}：把中心节点周围的图结构，用两种模板之一重排成定长的“节点序列”——
  上方 \emph{Neighborhood Detail Template}（邻域细节模板）把邻居展开成一棵计算树再按层遍历；
  下方 \emph{Hop-Field Overview Template}（跳数视野概览模板）把每一跳的邻居汇总成一个向量。
  \textbf{Step 2}：把节点序列经投影器（Projector，唯一可训练，标 “Tuned”）映射进 token 空间，
  连同提示词一起喂给\emph{冻结的} LLM（Freezed LLM，基座为 Vicuna-v1.5）。
  \textbf{右侧}：三种训练任务统一成问答格式——节点分类、链接预测、节点描述。}
  \label{fig:arch}
\end{figure}

\subsection{第 1 步：结构感知的“图$\to$节点序列”翻译（无参数）}

图分析的基本单元是\textbf{节点}。作者设计了\textbf{两种节点级模板}，从两个不同视角把“中心节点 + 它的邻域结构”
摊平成一条\textbf{定长}序列。注意：这一步\textbf{不含任何可学习参数}，只是按规则搬运，所以原始结构信息被完整保留、
无需“蒸馏”。

\textbf{模板一：邻域细节模板（Neighborhood Detail Template）}——看得细。
给定中心节点 $v$，先以它为根建一棵\textbf{形状固定的采样计算树}。为每一跳设一个采样数 $n_1,n_2,\dots$
（$n_i$ 是第 $i$ 跳要采几个邻居）。从 $v$ 的 1 跳邻居 $N_v^1$ 里随机选 $n_1$ 个；如果邻居不够 $n_1$ 个，
就用\textbf{占位节点 [pad]} 补齐到 $n_1$ 个。这样每个邻居集合大小都恒为 $n_1$。再对每个被选中的邻居，
递归地各采 $n_2$ 个邻居当它的孩子，不够同样用 [pad] 补。占位节点的孩子只能是占位节点。
最后对这棵树做\textbf{层序遍历（按层从左到右）}，就得到一条定长序列。

\begin{tip}{用图~\ref{fig:arch} 左上角的例子走一遍}
以 A 为根、两跳、每跳采 3 个。A 的一阶邻居是 \{B, C, D\}，放在第二层。B 只有邻居 \{A, G\}，
补成 \{A, G, [pad]\}；C、D 同理。最终得到一棵“3 叉满树”，层序遍历后序列是：
\texttt{A B C D A G [pad] A [pad] [pad] A E F}。\textbf{每个位置都唯一对应原图里的一个相对结构位置}——
这正是“结构感知”的含义。
\end{tip}

把序列里每个节点用文本编码器 $\phi$ 编码成向量后，作者还在每个位置\textbf{额外拼接一个拉普拉斯位置嵌入
（Laplacian Embedding）}$U_i$，进一步强化结构信息。记计算树的邻接矩阵为 $A_{tree}$，其拉普拉斯矩阵
做特征分解：
\begin{equation*}
  L = I - D^{-\frac{1}{2}} A_{tree} D^{-\frac{1}{2}} = U^{T}\Lambda U \tag{1}
\end{equation*}
其中 $D$ 是 $A_{tree}$ 的度矩阵，$U$ 就是这棵模板树的拉普拉斯（位置）嵌入。
\begin{tip}{为什么这个嵌入“只算一次”就够了}
因为采样数固定，\textbf{所有节点用的计算树形状都一模一样}，所以 $U$ 对所有图、所有节点都通用，
只需在一开始算一次。它的作用类似 Transformer 里的“位置编码”——告诉模型“序列里这个位置在树里处于什么结构角色”。
\end{tip}
于是第 $i$ 个位置的最终节点嵌入 $h_{v_i}$ 为（$\|$ 表示向量拼接）：
\begin{equation*}
  h_{v_i} =
  \begin{cases}
    \mathbf{0} \,\|\, U_i, & \text{若 } v_i = \text{[pad]（占位节点用全零向量）};\\
    \phi(x_{v_i}) \,\|\, U_i, & \text{否则（真实节点用其文本编码）}.
  \end{cases}\tag{2}
\end{equation*}
最终，中心节点与其结构信息被转成定长的节点嵌入序列 $h_{v_1}, h_{v_2}, \dots, h_{v_n}$。

\textbf{模板二：跳数视野概览模板（Hop-Field Overview Template）}——看得广。
它不为每个邻居单独留一个向量，而是\textbf{把每一“跳”的所有邻居汇总成一个向量}。对中心节点 $v$，
第 $i$ 跳的嵌入 $h_v^{i}$ 通过\textbf{无参数的消息传递}（就是对上一跳邻居向量求平均）得到：
\begin{equation*}
  h_v^{i} = \frac{1}{|N_v^{1}|}\sum_{v' \in N_v^{1}} h_{v'}^{\,i-1},
  \qquad\text{其中 } h_x^{0} = \phi(x_v).\tag{3}
\end{equation*}
这里 $h_x^0$ 是节点自身的文本编码（第 0 跳）；逐跳平均后，$h_v^{i}$ 就\textbf{浓缩了第 $i$ 跳邻域的整体信息}。
于是一串 $h_v^0, h_v^1, h_v^2,\dots$（每跳一个向量）就代表了中心节点和它的多跳结构。
\begin{tip}{两个模板怎么选？}
\textbf{细节模板}保留每个邻居，适合“需要看清具体某个邻居”的任务（比如判断商品是不是“任天堂 Switch”那一类）；
\textbf{概览模板}用一个向量概括一整跳，视野更大但牺牲了细节，适合“看大局”的任务（比如看一篇论文周围
大体上都是哪些领域）。论文里把用细节模板的版本记作 \textbf{LLaGA-ND}，用概览模板的记作 \textbf{LLaGA-HO}。
\end{tip}

\subsection{第 2 步：用投影器把“节点序列”对齐进 token 空间}

要让 LLM 看懂图输入，必须把\textbf{节点嵌入空间}和 LLM 的\textbf{输入 token 空间}对齐。做法是：用一个
专门校准的投影器 $f_\theta$（实现上就是一个简单的 MLP），把每个节点嵌入 $h_i$ 映射成一个 token 嵌入 $e_i$：
\begin{equation*}
  e_i = f_\theta(h_i).\tag{4}
\end{equation*}
于是节点嵌入序列 $h_1, h_2, \dots, h_n$ 就被转成对应的 token 嵌入序列 $e_1, e_2, \dots, e_n$。
\textbf{划重点：投影器的参数 $\theta$ 是整个训练过程中\emph{唯一}被更新的参数。}

\subsection{第 3 步：对齐微调（Alignment Tuning）——训练目标}

作者用前面那三种任务（节点分类、链接预测、节点描述）一起来训练这个投影器。把这三种任务混在一起训，
能让投影器\textbf{从多个角度理解图}，练成一个“图嵌入$\leftrightarrow$token 嵌入”的通用翻译器，并顺带获得
生成节点解释的能力。

训练时所有数据都写成\textbf{聊天问答格式}（基座 LLM 用 Vicuna，按其规范设好系统提示）。处理输入时，
把提示词里的普通文字正常 tokenize；唯独遇到占位符 \texttt{<node sequence>} 时，
就\textbf{用投影后的节点嵌入 $e_1,\dots,e_n$ 原位替换}它。训练目标很直接——\textbf{最大化生成正确答案的概率}：
\begin{equation*}
  \underset{\theta}{\text{maximize}}\ \ p\big(X_{\text{answer}} \mid X_{\text{graph}}, X_{\text{question}}, X_{\text{system}}\big).\tag{5}
\end{equation*}
其中 $X_{\text{graph}}$ 是（投影后的）图序列、$X_{\text{question}}$ 是问题、$X_{\text{system}}$ 是系统提示、
$X_{\text{answer}}$ 是标准答案。注意优化变量只有 $\theta$（投影器），\textbf{LLM 本体全程冻结}。

\begin{tip}{一句话抓住整套方法}
LLaGA = 把节点邻域\textbf{规则化重排}成“结构感知序列”（无参数，含位置嵌入）$\to$ 经\textbf{一个小投影器}
变成 LLM 词向量空间里的“外语词” $\to$ 连同问题一起喂给\textbf{冻结的} LLM 生成答案。全程只训投影器，
所以又轻、又快、又不破坏 LLM 的通用能力。
\end{tip}

\section{七、实验怎么做、结果怎么读}

\textbf{实验设置}：基座 LLM 默认用 \textbf{Vicuna-7B-v1.5-16K}；文本编码默认用 \textbf{SimTeG}；学习率 2e-5、
批大小 16、只训 1 个 epoch（最小的数据集 Cora 复制 3 倍以补数据量）。细节模板采样 2 跳、每跳 10 个邻居；
概览模板用 4 跳嵌入。对比的基线分三类：GNN（GCN、GraphSAGE、GAT、SGC、SAGN）、图 Transformer（NodeFormer）、
以及通用大模型 GPT-3.5；还和同期工作 GraphGPT 比。

\textbf{四种训练设定}（一个比一个“通才”，越靠后越难，越能体现 LLaGA 价值）：
\begin{itemize}[leftmargin=2em]
  \item \textbf{Single Focus}：单数据集、单任务（最专精）。
  \item \textbf{Task Expert}：跨所有数据集、但只做单一任务。
  \item \textbf{Classification Expert}：跨所有数据集，同时做节点分类和链接预测两项。
  \item \textbf{General Model}：跨所有数据集，同时做\textbf{三项任务}（含节点描述）——最“全才”。
\end{itemize}
基线 GNN 做多任务时，靠“共享主干 + 多个任务专用分类头”；而 LLaGA 始终只用\textbf{一个投影器}通吃。

\textbf{主结果（节选自原文 Table 1）}。下表挑了有代表性的几行：在最简单的 Single Focus 设定里，
LLaGA 已能与最强 GNN 平起平坐甚至更好；越到“多任务”设定，差距越明显。

\begin{table}[H]\centering
\caption{节点分类与链接预测准确率（\%）对比（节选自原文 Table 1；越高越好，加粗为该设定下最优）。}
\begin{tabular}{llcccc}
\toprule
设定 & 模型 & Arxiv & Products & Pubmed & Cora \\
\midrule
\multirow{2}{*}{Single Focus（分类）}
 & GraphSAGE     & 76.29 & 82.87 & 94.87 & 88.89 \\
 & \textbf{LLaGA-HO-7B} & \textbf{76.66} & \textbf{84.67} & 95.03 & \textbf{89.22} \\
\midrule
\multirow{2}{*}{Classification Expert（分类）}
 & GCN           & 70.95 & 80.02 & 89.00 & 82.77 \\
 & \textbf{LLaGA-HO-7B} & \textbf{75.99} & \textbf{83.32} & \textbf{94.80} & \textbf{89.30} \\
\midrule
\multirow{2}{*}{Classification Expert（链接预测）}
 & GCN           & 87.69 & 92.88 & 72.28 & 78.35 \\
 & \textbf{LLaGA-HO-7B} & \textbf{94.30} & \textbf{96.05} & \textbf{88.64} & \textbf{88.53} \\
\bottomrule
\end{tabular}
\end{table}

\textbf{怎么读这张表}：作者总结出三点观察。\textbf{观察 1}——LLaGA 在所有设定下都优于基线，尤其在多任务场景。
\textbf{观察 2}——很多基线一进多任务场景就\textbf{大幅掉分}（比如上表里 GCN 在 Classification Expert 的
Pubmed 链接预测只有 72.28），而 LLaGA \textbf{几乎不掉甚至还升}，说明它擅长跨数据/任务提取共性。
\textbf{观察 3}——两个模板各有所长：细节模板擅长需要具体邻居信息的任务，概览模板擅长需要大视野的任务。

\textbf{对比同期工作 GraphGPT（源自原文 Table 2）}：LLaGA 最通用的模型一口气训了 12 个任务，
而 GraphGPT 最通用的模型只训了 3 个任务；即便如此，LLaGA 在这三项上仍\textbf{全面胜出}（例如 Pubmed 链接预测
89.81 对 GraphGPT 的 80.26）。

\textbf{节点描述与可解释性（RQ2，源自原文 Table 4）}：LLaGA 生成的描述与真值的 Sbert 相似度\textbf{远高于}
“两条随机样本”的 Base 基准值（例如 Pubmed 上 0.69 对 0.49），且从描述里反推类别的准确率也很高。
原文 Table 3 还展示了实例：即便偶尔预测的标签和真值不同，它给出的解释也\textbf{合情合理}（比如把一篇抑郁症检测论文
判成“社交与信息网络”而非“机器学习”，但理由站得住脚）。

\textbf{零样本迁移（RQ3，源自原文 Table 5 / Table 8）}：训练时没见过的数据集，测试时直接上。
这里 LLaGA 的优势最抢眼：

\begin{table}[H]\centering
\caption{链接预测的零样本迁移准确率（\%）（数值源自原文 Table 5；越高越好）。}
\begin{tabular}{llc}
\toprule
训练 $\to$ 测试 & 模型 & 准确率 \\
\midrule
\multirow{2}{*}{Arxiv+Pubmed $\to$ Cora（同领域）}
 & GraphGPT-7B   & 50.74 \\
 & \textbf{LLaGA-HO-7B} & \textbf{87.35} \\
\midrule
\multirow{2}{*}{Arxiv+Pubmed+Cora $\to$ Products（跨领域）}
 & GraphGPT-7B   & 50.74 \\
 & \textbf{LLaGA-HO-7B} & \textbf{92.99} \\
\bottomrule
\end{tabular}
\end{table}

注意第二行是\textbf{跨领域}迁移（训练全是引文网络，测试却是电商网络），LLaGA 仍达到 92.99\%，
说明它学到的“图$\leftrightarrow$token 对齐”能\textbf{迁移到差异很大的新领域}。节点分类的零样本同样大幅领先
GraphGPT（原文 Table 8，例如 Arxiv+Pubmed$\to$Cora 在“节点嵌入+文本属性”提示下达到 59.59\%，
GraphGPT 仅 44.65\%）。

\textbf{模板消融（RQ4，源自原文 Table 6）}：去掉模板、只用中心节点自己的嵌入做预测，性能明显下降——
尤其在\textbf{高度依赖结构信息}的链接预测上掉得最多（例如 Pubmed 从约 92.36 掉到 78.19）。
这证明\textbf{两个结构模板确实把邻域结构信息有效编码进去了}。

\textbf{鲁棒性补充（附录 C / D）}：换文本编码器（SBERT、RoBERTa）或换基座 LLM（LLaMA2-7B、甚至更轻的 OPT-2.7B），
LLaGA 都\textbf{稳定优于 GNN 基线}，说明方法对“编码器/基座”的选择不挑食。

\section{八、附录里值得一看的补充}
\begin{itemize}[leftmargin=2em]
  \item \textbf{数据划分}（附录 A）：节点级任务用标准 train/val/test 划分（Arxiv 6:2:3、Products 8:2:90、
        Pubmed 与 Cora 均 6:2:2）；链接预测的训练/测试节点对分别从对应集合里采样。
  \item \textbf{节点分类零样本更难}（附录 B）：因为不同任务\emph{标签集不同}，比链接预测更具挑战；作者借助
        “节点描述任务”先建立图与语义空间的通用对齐，再迁移到分类，效果仍超 GraphGPT。
  \item \textbf{换文本编码器}（附录 C，Table 9）：SBERT / RoBERTa 都能用，LLaGA 始终领先 GNN。
  \item \textbf{换基座 LLM}（附录 D，Table 10）：Vicuna-7B / LLaMA2-7B / OPT-2.7B 都行，连 2.7B 的小模型也有效。
  \item \textbf{方差信息}（附录 E，Table 11）：在 Cora、Pubmed 上重复 5 次，结果方差很小（如 Cora 分类 88.86$\pm$0.78），说明稳定。
\end{itemize}

\section{九、小白知识点卡片（速查）}
\begin{itemize}[leftmargin=2em]
  \item \textbf{文本属性图 / TAG}：节点带文字的图（引文网络、电商网络等）。
  \item \textbf{GNN（图神经网络）}：靠“消息传递+聚合”把节点编码成向量的网络；强但不擅长多任务。
  \item \textbf{token 嵌入空间 / 词向量空间}：LLM 把每个词查成向量后所在的空间。
  \item \textbf{投影器（projector）}：一个小 MLP，把图节点向量“搬进”词向量空间当“外语词”；本文唯一可训练部件。
  \item \textbf{结构感知序列}：把中心节点的邻域按模板规则摊平成的定长序列，每个位置对应一个固定的结构角色。
  \item \textbf{邻域细节模板 / 跳数视野概览模板}：两种重排方式，一个“看得细”、一个“看得广”（对应 LLaGA-ND / LLaGA-HO）。
  \item \textbf{拉普拉斯位置嵌入}：给序列每个位置加的“结构位置编码”，因模板形状固定故只算一次。
  \item \textbf{占位节点 [pad]}：邻居不够时补的空节点，用全零向量表示。
  \item \textbf{冻结 LLM（frozen LLM）}：训练时完全不改 LLM 参数，保住其通用能力。
  \item \textbf{零样本（zero-shot）}：在训练时\emph{没见过}的数据集/任务上直接评测，考验迁移能力。
  \item \textbf{Vicuna / SimTeG / SBERT}：分别是本文用的基座 LLM、默认文本编码器、可选文本编码器。
\end{itemize}

\section{十、一句话总结与定位}

\textbf{一句话}：LLaGA 把“一个节点 + 它的邻域结构”用\emph{无参数的模板}重排成“结构感知 token 序列”，
再用\emph{一个小投影器}把它映射进 LLM 的词向量空间当“外语词”，\emph{冻结整个 LLM、只训这个投影器}，
就让\emph{一个统一模型}同时做好节点分类、链接预测、节点描述，还能零样本迁移到全新的数据集和领域。

\textbf{在“图 $\times$ 后训练”谱系里的位置}：LLaGA 是\textbf{“表示对齐 / 投影器路线”}的代表作——
它不改 LLM、不靠预训练图编码器、也不把图“说成”啰嗦的自然语言，而是把图\textbf{直接编码成序列后投影对齐}。
与之相对：同系列里 G-Retriever 走“图上检索（GraphRAG）+ 软提示”、InstructGraph/GraphWiz 走“指令微调 + DPO”、
G1 走“强化学习”。你可以对照着看它们\emph{各自把“图的信息”塞进 LLM 的方式有何不同}——LLaGA 的特色就是
\textbf{“轻投影器对齐 + 冻结基座 + 一模型多任务 + 强零样本”}。

\textbf{动手复现}：官方代码 \url{https://github.com/VITA-Group/LLaGA}（含两种模板的构造、投影器训练、
以及 Vicuna 基座的统一问答训练流程），是入门“图 $\times$ LLM 表示对齐”很好的练手项目。

\end{document}
